% Case 1: reaction at the roller as the redundant force
\begin{tikzpicture}[scale=0.7,font=\small]
  \begin{scope}
    \pic at (0,0) {wallL}; \draw[member] (0,0) -- (6,0);
    \node[above right] at (0,0.05) {$a$}; \node[above right] at (3,0.05) {$b$}; \node[above left] at (6,0.05) {$C$};
    \draw[load] (3,1.3) -- (3,0.1) node[pos=0,above] {$P$};
    \draw[load] (6,-1.2) -- (6,-0.1) node[pos=0,below] {$R$};
  \end{scope}
  \begin{scope}[xshift=8cm]
    \draw[thinmember] (0,0) -- (3.5,0); \draw[thinmember] (6.5,0) -- (10,0);
    \node[above right] at (0.05,0) {$a$}; \node[above left] at (3.5,0) {$b$}; \node[above right] at (6.55,0) {$b$}; \node[above left] at (10,0) {$C$};
    \draw[->,thick] (-0.15,0.45) arc (110:-110:0.45); \node[above] at (-0.2,0.6) {$\frac{PL}{2}-RL$};
    \draw[->,thick] (3.65,0.45) arc (70:290:0.45); \node[above] at (3.7,0.6) {$\frac{RL}{2}$};
    \draw[->,thick] (6.35,0.45) arc (110:-110:0.45); \node[above] at (6.4,0.6) {$-\frac{RL}{2}$};
    \draw[->,thick] (10.15,0.45) arc (70:290:0.45); \node[above] at (10.2,0.6) {$0$};
  \end{scope}
\end{tikzpicture}
